The Gravitational-Wave Optical Transient Observer (GOTO) aims to carry out
searches of optical counterparts of any Gravitational Wave (GW) sources within
an hour of their detection by the LIGO and VIRGO experiments. The survey needs
to cover a big area on the sky in a short time to search all possible regions
constrained by the GW experiments before the optical burst signal, if any,
fades away. The GOTO telescope array is composed of 4 identical reflective
telescopes, each with 40-cm diameter aperture and f-ratio of 2.5, working
together to image large regions of the sky. For simplicity, we assume that the
telescopes' fields-of-view (FoV) do not overlap with one another.

\begin{description}
\item[a)] Calculate the projected angular size per mm at the focal plane, i.e.\
  plate scale, of each telescope. \hfill[6]
\item[b)] If the zero-point magnitude (i.e.\ the magnitude at which the count
  rate detected by the detector is 1 count per second) of the telescope system
  is 18.5\,mag, calculate the minimum time needed to reach 21\,mag at
  Signal-to-Noise Ratio (SNR) $= 5$ for a point source. We first assume that
  the noise is dominated by both the Read-Out Noise (RON) at 10 counts/pixel
  and the CCD dark (thermal) noise (DN) rate of 1 count/pix/minute. The CCDs
  used with the GOTO have a 6-micron pixel size and gain (conversion factor
  between photo-electron and data count) of 1. The typical seeing at the
  observatory site is around 1.0\,arcsec. \hfill[8]\\
  The Signal-to-Noise Ratio is defined as
  \[ \mathrm{SNR} \equiv
     \frac{\text{Total Source Count}}{\sqrt{\sum_i \mathrm{Noise}_i^2}}
     = \frac{\text{Total Source Count}}
            {\sqrt{\sigma_{\mathrm{RON}}^2 + \sigma_{\mathrm{DN}}^2 + \dots}},
  \]
  \[ \sigma_{\mathrm{RON}} = \sqrt{N_{\mathrm{pix}} \cdot \mathrm{RON}^2},
     \qquad
     \sigma_{\mathrm{DN}} = \sqrt{N_{\mathrm{pix}} \cdot \mathrm{DN} \cdot t},
  \]
  where $t$ is the exposure time.
\item[c)] Normally when the exposure time is long and the source count is high
  then Poisson noise from the source is also significant. Determine the
  relation between SNR and exposure time in the case that the noise is
  dominated by Poisson noise of the source. Recalculate the minimum exposure
  time required to reach 21\,mag with SNR $= 5$ from part b) if Poisson noise
  is also taken into consideration. The Poisson noise (standard deviation) of
  the source is given by
  $\sigma_{\mathrm{source}} = \sqrt{\text{Source Count}}$. In reality, there is
  also the sky background which can be important source of Poisson noise. For
  our purpose here, please ignore any sky background in the calculation.
  \hfill[6]
\item[d)] The typical localisation uncertainty of the GW detector is about 100
  square-degrees and we would like to cover the entire possible location of any
  candidate within an hour after the GW is detected. Estimate the minimum side
  length of the square CCD needed for each telescope in terms of the number of
  pixels. You may assume that the time taken for the CCD read-out and the
  pointing change are negligible. \hfill[5]
\end{description}
